\scanpage{19}

ils correspondent aux topologies \inserted{$T'$} sur $C$ \deleted{\unsure{plus}}\inserted{plus} fines que la top. donnée, par $R'\longleftrightarrow\widetilde{(C,T')}$.

Remarquons que pour un sous topos donné \inserted{$R'$}, on peut tjrs le décrire à l'aide d'un \uline{petit} ens. $\Sigma\subset\mathrm{Fl}\,R$. En effet, soit $C\subset R$ \inserted{petite et} génératrice \deleted{\unsure{de}}. Alors les $\Delta$-flèches $\in\Delta$ $R\longrightarrow\unsure{E|R'|\bar R}$ qui se factorisent par $R'$ sont celles qui transforment familles \deleted{couvrantes} $S_i\to S$ \inserted{dans $C$ couvrantes pour la topologie $T'$ définissant $R'$} \deleted{\unsure{de $R$}} en familles couvrantes [ou encore telles que pour $\forall S\in\mathrm{Ob}\,R$ et $S'\hookrightarrow S$ monomorphisme de \deleted{\unsure{$E$}} $R$, $S'\to S$ est transformé en un \uline{iso}].

\uline{Conséquence} Supposons $R$ ``algébrique'' i.e. $T_R$ \inserted{théorie} algébrique, alors $T_{R'}\subset T_R$ (ss-\deleted{\unsure{catég|catégorie}}théorie ``pleine'') est algébrique, déduite de $T_{R'}$ par un petit ensemble d'axiomes.

D'ailleurs, si on a deux théorie représentables $T_R,T_{R'}$ ($R,R'$ des topos) et un foncteur $T_{R'}\unsure{-|\to} T_R$ qui est un plongement en théories, i.e. tel que \deleted{$\underline{\mathrm{Hom}}_\Delta$} $T_{R'}(E)-T_R(E)$ pl. fid. pour tt $E\in\mathrm{Ob}\,\Delta$, \deleted{\unsure{i.e.} $\underline{\mathrm{Hom}}_\Delta($} alors la flèche correspondante $R\to R'$ de $\Delta$ est telle que $\forall E\in\mathrm{Ob}\,\Delta$, $\underline{\mathrm{Hom}}_\Delta(R',E)-\underline{\mathrm{Hom}}_\Delta(R,E)$ est pl. fidèle. Or ceci caractérise les morphismes de topos $R'-R$ qui sont des plongements.

4) \uline{\unsure{Théorèmes|Théorème}} a) Un $R\in\mathrm{Ob}\,\Delta$ est algébrique, i.e. représente une théorie algébrique, ssi $R$ est un topos, i.e. admet une petite sous-catégorie génératrice. .

b) Soit $I$ un cat discrète, $R\in\mathrm{Ob}\,\Delta$, $I\overset b{-} R$ un foncteur. Alors \unsure{$*$} $b$ est basique (i.e. ...) \deleted{\unsure{désigne|disque}}\deleted{\unsure{ssi} $b(I)$ est génératrice} ssi la sous-catégorie pleine de $R$ engendrée par produits finis d'objets de $b(I)$ est strictement génératrice.

\uline{Dém} \deleted{\unsure{Évidemment b) $\Rightarrow$ a)}} \inserted{Évidemment b) $\Rightarrow$ a, donc} prouvons \deleted{\unsure{d'abord}} b).

\uline{Suffisance} On va construire la théorie alg. \inserted{basique} représentée par $R,b$ en plusieurs pas.

% archive sheet number: 18
